% ============================================================
\section{Notation and Toolkit}
\label{sec:orientation}
% ============================================================

\subsection{The Problem: Finite Observability}

When an observer has finite resolution, not all structure is representable.
A projection $\CCPi{h}$ captures ``what can be seen at resolution $h$.''
The fundamental question:
\begin{center}
\emph{What is lost, and what defects arise, when operators are applied under finite observability?}
\end{center}

\subsection{The Interface Objects}

\begin{center}
\renewcommand{\arraystretch}{1.3}
\begin{tabular}{@{}lll@{}}
\textbf{Object} & \textbf{Notation} & \textbf{Role} \\
\hline
Additive state space & $\cH$ & Ambient space for structure \\
Horizon projection & $\CCPi{h}$ & What the observer can represent \\
Shadow projection & $\CCPibar{h} = \id - \CCPi{h}$ & What escapes observation \\
Field-induced projection & $\CCSpProj{\CCField}{S}$ & Spectral cut on field $\CCField$ \\
Observed operator & $\CCObs{A}{h} = \CCRestr{A}{h}$ & The operator as seen \\
Leakage operator & $\CCLeak{A}{h} = \CCLeakExpr{A}{h}$ & Observable $\to$ shadow flow \\
Defect operator & $\CCDefect{A}{h} = \CCDefectExpr{A}{h}$ & Round-trip through shadow \\
\end{tabular}
\end{center}

\subsection{The Tower Calculus}

The \emph{tower calculus} is the algebraic spine of Coherence Calculus.
It applies to \emph{any} operator on a discrete horizon tower; no nilpotency
assumption is required.

\smallskip\noindent
\textbf{Calculus completeness.}
The framework is internally complete in exactly three components:
\begin{enumerate}[label=(\roman*), itemsep=0.1em, leftmargin=2em]
\item \textbf{Tower calculus}: derivative $\partial$, antiderivative $\CCInt{\cdot}$, FTC (Section~\ref{subsec:tower-antiderivative}; Theorem~\ref{thm:ftc-tower}).
\item \textbf{Defect grammar}: chain/product rules for $\CCTau{h}{F}{x}$ (Section~\ref{sec:defect-rules}).
\item \textbf{Elimination/resummation}: split transform, Schur correction $\CCSigma{h}{z}$ (Section~\ref{sec:resolvent-defect}).
\end{enumerate}
These are not separate theories; they form a unified grammar for manipulating horizon-restricted operators.
The completion-and-selection contract theorem (Section~\ref{sec:constraint-compiler}) makes this composition explicit: defect and elimination produce a bound \(D\), and admissibility selects a minimal completion together with its forcing or residual data.

\begin{center}
\renewcommand{\arraystretch}{1.3}
\begin{tabular}{@{}lll@{}}
\textbf{Calculus Role} & \textbf{Object} & \textbf{Identity} \\
\hline
Derivative & $\CCDer{A} = \CCComm{\CCIndexOp}{A}$ & Leibniz rule, block formula \\
Integral & $\CCInt{B}_{hk} = \frac{1}{h-k} B_{hk}$ & Right-inverse on off-diagonal \\
FTC analog & $A = \CCDiag{A} + \CCInt{\CCDer{A}}$ & Decomposition \\
IBP & $\Tr(\CCDer{A} B) = -\Tr(A \CCDer{B})$ & Trace pairing \\
\hline
\multicolumn{3}{@{}l@{}}{\emph{Elimination (any projection $P$):}} \\
Split transform & $\CCSplit{P}{A}{z} = (\CCAeff{A}{P}, \CCSigmaP{P}{z})$ & Effective + residual \\
Inverse design & $\CCDesign{P} = \|\CCSigmaP{P}{z} - \Sigma^\star\|^2$ & Existence on $\CCGrass{r}$ \\
\hline
\multicolumn{3}{@{}l@{}}{\emph{Solver layer (when graded complex is exact):}} \\
Hodge Laplacian & $\CCLap{k} = D_{k-1}D_{k-1}^* + D_k^* D_k$ & Invertible if $\CCCoh{k}(D)=0$ \\
Canonical integration & $\CCHom{k} = D_{k-1}^* \CCGreen{k}$ & Homotopy: $DK + KD = I$ \\
\end{tabular}
\end{center}

% ------------------------------------------------------------
\begin{tcolorbox}[colback=white, colframe=black!70, title={Certificate Checklist}]
\label{box:certificate-checklist}
Three certifications gate the main capabilities of Coherence Calculus:

\smallskip
\textbf{\CertA\ Perfect Coherence (Zero Residual):}
Check $\CCBlockPQ{A} = 0$ or $\CCBlockQP{A} = 0$.\\
\emph{If satisfied:} No Schur correction needed; observed dynamics are exact.\\
\emph{Reference:} Theorem~\ref{thm:block-diagonal-sufficient} (Section~\ref{sec:targets}).

\smallskip
\textbf{\CertB\ Solver Exactness:}
Check $\CCCoh{k}(D) = 0$ (cohomology vanishes at degree $k$).\\
\emph{If satisfied:} Green operator $\CCGreen{k} = \CCLap{k}^{-1}$ and homotopy $\CCHom{k}$ exist.\\
\emph{Reference:} Definition~\ref{def:exactness}, Theorem~\ref{thm:homotopy-identity} (Section~\ref{subsec:solver-layer}).

\smallskip
\textbf{\CertC\ Holographic Saturation:}
Check $\Ker(R_k) = \Img(D_{k-1})$ (relative exactness).\\
\emph{If satisfied:} Zero boundary data implies pure gauge (bulk state is exact).\\
\emph{Reference:} Definition~\ref{def:holographic-saturation} (Section~\ref{subsec:solver-layer}).

\smallskip
\emph{The admissible-completion layer (Section~\ref{sec:constraint-compiler}) uses these certificates as gates in the selection pipeline.}
\end{tcolorbox}
% ------------------------------------------------------------

% ------------------------------------------------------------
\begin{tcolorbox}[colback=white, colframe=black!70, title={Taxonomy of Operators}]
\label{box:taxonomy}
\textbf{Problem data} (given by the application):
\begin{itemize}[itemsep=0.2em, leftmargin=1.5em]
\item \textbf{Problem differential} $d$ (or graded $D_k: H^k \to H^{k+1}$): the closure operator of interest, satisfying $d^2 = 0$.
\end{itemize}

\smallskip
\textbf{Horizon-restricted operators} (derived from $d$ and $\CCPi{h}$):
\begin{itemize}[itemsep=0.2em, leftmargin=1.5em]
\item \textbf{Observed differential} $\CCObs{d}{h} = \CCRestr{d}{h}$: what the observer sees when applying $d$.
\item \textbf{Leakage operator} $\CCLeak{d}{h} = \CCLeakExpr{d}{h}$: the part of $d$ that sends observable states into shadow.
\item \textbf{Defect operator} $\CCDefect{d}{h} = \CCDefectExpr{d}{h}$: round-trip through shadow; measures $\CCObs{d}{h}^2 = -\CCDefect{d}{h}$.
\end{itemize}

\smallskip
\textbf{Tower calculus operators} (the calculus machinery):
\begin{itemize}[itemsep=0.2em, leftmargin=1.5em]
\item \textbf{Tower derivative} $\CCDer{A} = \CCComm{\CCIndexOp}{A}$: \emph{the} calculus derivative; acts on operators, extracts cross-scale transport with scale-separation weights. Satisfies Leibniz rule.
\item \textbf{Tower antiderivative} $\CCInt{B}$: right-inverse of $\partial$ on off-diagonal blocks; $\CCBlock{(\CCInt{B})}{h}{k} = \frac{1}{h-k} \CCBlock{B}{h}{k}$ for $h \ne k$. Recovers off-diagonal structure from derivative data.
\end{itemize}

\smallskip
\textbf{Discrete update operators} (NOT calculus derivatives):
\begin{itemize}[itemsep=0.2em, leftmargin=1.5em]
\item \textbf{Horizon flow} $\CCFlow{A}{h} = A_{h+1} - A_h$: finite difference of observed operators across adjacent horizons. A discrete update, not a commutator derivative.
\item \textbf{Increment projector} $\CCDelta{h} = \CCPi{h} - \CCPi{h-1}$: projects onto the ``layer'' added at resolution $h$.
\end{itemize}

\smallskip
\textbf{Key disambiguation:}
$\partial$ and $\CCInt{\cdot}$ are the calculus operators (derivative/antiderivative on the operator algebra).
$d$ is problem data (the differential whose behavior under projection we study).
$\nabla_h$ is a discrete update, not a derivative.
\end{tcolorbox}
% ------------------------------------------------------------

\subsection{The Horizon Fundamental Theorem}

The \emph{Horizon Fundamental Theorem} applies specifically to \emph{graded
differentials} ($d^2 = 0$) under horizon projection. It is an exact identity
from block algebra. Its role is organizational: it gives a compact entry point
to the defect grammar developed later in this part.

\begin{tcolorbox}[colback=black!3, colframe=black!50, title={Horizon Fundamental Theorem (Preview)}]
\textbf{HFT-1 (Defect Identity):} If $d^2 = 0$, then $\CCObs{d}{h}^2 = -\CCDefect{d}{h}$.

\smallskip
\textbf{HFT-2 (Telescoping):} Total deviation = sum of orthogonal layer increments.

\smallskip
\emph{These identities organize the chain rules, product rules, and elimination patterns of the later chapters of Part~II.}
\end{tcolorbox}

\subsection{Running Example: The 4-Vertex Path}
\label{subsec:hello-world}

\begin{example}[Running example used throughout]
\label{ex:hello-world}\lean{CoherenceCalculus.pathGraph4, CoherenceCalculus.blockAverageProj4}
Let $\cH = \bbR^4$ with standard basis $e_1, e_2, e_3, e_4$, and let $P$ be
the block-averaging projection defined below.

\textbf{Differential} (degree-0$\to$1 map in a graded complex; grading suppressed):
\[
(Dx)_1 = x_2 - x_1, \quad (Dx)_2 = x_3 - x_2, \quad (Dx)_3 = x_4 - x_3.
\]

\textbf{Horizon projection} (block averaging $\{1,2\}$ and $\{3,4\}$):
\[
(P x)_1 = (P x)_2 = \tfrac{x_1+x_2}{2}, \quad
(P x)_3 = (P x)_4 = \tfrac{x_3+x_4}{2}.
\]

\textbf{Observable operator:} $PDP$ sees only inter-block differences.

\textbf{Leakage:} $(I-P)DPx$ captures within-block differences that escape observation.

\textbf{Defect energy:} For $x = (0,1,1,0)$:
$Dx = (1,0,-1)$, $Px = (\frac{1}{2}, \frac{1}{2}, \frac{1}{2}, \frac{1}{2})$, $DPx = (0,0,0)$.
Thus $\|Dx\|^2 - \|DPx\|^2 = 2$.

\emph{This example will be referenced as ``the 4-vertex path'' throughout.}
\end{example}

% --- Horizon Tower Diagram ---
\begin{figure}[ht]
\centering
\resizebox{0.82\textwidth}{!}{%
\begin{tikzpicture}[node distance=1.5cm]
  \node[font=\small\bfseries] at (0,1.2) {Nested subspaces};
  \node[font=\small\bfseries] at (6.2,1.2) {Increment projectors};

  \node[ccbox] (H) at (0,0) {$\cH$};
  \node[ccbox, below=of H] (H2) {$\CCHobs{2}$};
  \node[ccbox, below=of H2] (H1) {$\CCHobs{1}$};
  \node[ccbox, below=of H1] (H0) {$\CCHobs{0}$};

  \draw[ccarrow] (H) -- node[right, cclabel] {$\CCPi{2}$} (H2);
  \draw[ccarrow] (H2) -- node[right, cclabel] {$\CCPi{1}$} (H1);
  \draw[ccarrow] (H1) -- node[right, cclabel] {$\CCPi{0}$} (H0);

  \node[ccband, minimum width=1.5cm, minimum height=0.8cm] (D2) at (4.8,-1.5) {$\CCDelta{2}$};
  \node[ccband, minimum width=1.5cm, minimum height=0.8cm] (D1) at (4.8,-3.0) {$\CCDelta{1}$};
  \node[ccband, minimum width=1.5cm, minimum height=0.8cm] (D0) at (4.8,-4.5) {$\CCDelta{0}$};

  \node[right=0.55cm of D2, font=\scriptsize] {$= \CCPi{2} - \CCPi{1}$};
  \node[right=0.55cm of D1, font=\scriptsize] {$= \CCPi{1} - \CCPi{0}$};
  \node[right=0.55cm of D0, font=\scriptsize] {$= \CCPi{0}$};
\end{tikzpicture}%
}
\caption{Horizon tower with nested subspaces (left) and increment projectors as ``bands'' (right).}
\label{fig:horizon-tower}
\end{figure}

% --- Defect Round-Trip Diagram ---
\begin{figure}[ht]
\centering
\begin{tikzpicture}[node distance=3cm, scale=1.1, every node/.style={scale=1.1}]
  % Observable and shadow
  \node[ccbox, minimum width=2.2cm] (obs) {$\cH_{\mathrm{obs}}$};
  \node[ccbox, minimum width=2.2cm, right=of obs] (sh) {$\cH_{\mathrm{sh}}$};

  % Arrows forming round trip
  \draw[ccarrow, bend left=25] (obs.north east) to node[above, cclabel] {$\CCPibar{h} d$} (sh.north west);
  \draw[ccarrow, bend left=25] (sh.south west) to node[below, cclabel] {$\CCPi{h} d$} (obs.south east);

  % Self-loop for observed operator
  \draw[ccarrow] (obs.south) arc[start angle=-90, end angle=-270, radius=0.55cm]
    node[left, pos=0.5, cclabel] {$\CCObs{d}{h}$};

  % Annotation
  \node[below=1.6cm of obs, font=\small, align=center] {
    Round trip: $\CCDefect{d}{h} = \CCPi{h} d \CCPibar{h} d \CCPi{h}$
  };
\end{tikzpicture}
\caption{Defect as ``round trip through shadow'': leakage out ($\CCPibar{h} d$) and return ($\CCPi{h} d$).}
\label{fig:defect-roundtrip}
\end{figure}
